%% Kaynak: tez-latex-v2/frontmatter/abstract.tex (15.11.2026)
\begin{summary}
The performance of an energy storage system depends as much on the efficiency, reliability, volume and cost of the power conversion interface between the battery, the DC bus and the grid as on the battery pack itself. Grid-connected photovoltaic and storage converters operate at rated power only for limited periods and spend a considerable share of their operating life at partial load, where fixed driving, switching and magnetic losses grow relative to the transferred power and the current waveform becomes more sensitive to switch nonlinearities, dead time and grid distortion. The low-power region is therefore the operating range in which the trade-off between efficiency and current total harmonic distortion (THD) becomes most visible.

The three-level T-type inverter is an attractive grid interface owing to its high efficiency and power density at rated load. In the low-power region, however, its behaviour is governed by a different mechanism. The bidirectional neutral branch that produces the zero-voltage state keeps conducting after the inductor current reaches zero, so the phase leg behaves like a synchronous buck converter and lets the current reverse within a switching period. The reverse circulating current carries no useful power, yet it increases the RMS current and the conduction losses of the semiconductors and filter elements. Blocking reverse current is a known principle for synchronous buck converters, but applying it to the zero state of a grid-connected three-level T-type inverter changes the current path, the pole voltage, the modulation integrity and the harmonic content at the same time, and is therefore a topology-specific problem.

The thesis addresses the problem with three original contributions built on the T-type topology. The first two operate at the phase-leg scale, the third at the system scale. (1) A neutral-switch suppression method blocks the neutral-switch command that sustains reverse conduction. (2) A Sin\textsuperscript{3}-based feedforward compensation reduces the third-harmonic penalty of the suppression without altering the stability of the current loop. (3) The GT-type reconfigurable energy storage architecture re-tasks the same T-type semiconductors as a grid inverter, a Vienna-type power factor correction (PFC) stage and an active battery management stage. Four research questions are answered: the mechanism of reverse circulation, the efficiency--harmonic trade-off of suppression, the harmonic effect and loop interaction of the Sin\textsuperscript{3} term, and the feasibility of safe hardware re-tasking.

\textbf{System, model and method.} The prototype is a three-phase, four-wire T-type inverter whose neutral is tied to the DC-bus midpoint; simulations and experiments are carried out on one phase leg with an LCL filter ($L_1=600~\mu$H, $C_f=8~\mu$F, grid-side equivalent $L_2$), a switching frequency of 10~kHz and a dead time of 1~$\mu$s read from the PWM timer registers. The switching states, the four basic current paths and the synchronous-buck equivalent of the phase leg are drawn on a common circuit base, and the reverse-current threshold is derived from the condition that the average inductor current is smaller than half of the ripple. Four switched Simscape models are used: classical complementary switching (S1), neutral-switch suppression (S2), suppression with Sin\textsuperscript{3} feedforward (S3), and a closed-loop current-feedback model (S4) built in the thesis on top of S3. The gate logic read from the models confirms that the only difference between S1 and S2 is the opposite neutral-switch command in the zero state. The simulation models were calibrated to the experimental operating conditions, accounting for the equivalent electrical characteristics and losses of the experimental setup. All circuit, semiconductor, modulation, control and solver parameters are presented in a single table with their values, units and sources.

\textbf{Calibration against the experiment.} In the simulations, the equivalent grid and variac impedances were estimated from the experimental results, and inductor core losses were modelled empirically. The modulation-reference amplitude and phase were adjusted to reproduce the experimental operating points.

\textbf{Low-power results.} At a phase current of about 2.3~A, 250--280~W and a grid voltage of about 120~V, suppression raises the efficiency from 92.5\,\% to 94.2\,\% in the experiment and from 91.15\,\% to 93.93\,\% in the calibrated simulation. The ranking is the same in both sources. The price is a pronounced third-harmonic distortion: the grid-current THD rises from 15\,\% to 47\,\% in the experiment. With the Sin\textsuperscript{3} term added, the experimental THD falls to 14\,\%, back to the level of classical switching, and the efficiency reaches 94.5\,\%. The calibrated S3 model reproduces the AC current within 1.4\,\%, the efficiency within 0.25 points and, without THD being a calibration target, the current THD as 14.8\,\%.

\textbf{Sin\textsuperscript{3} compensation.} In the model the term is the cube of a unit-amplitude sine, added with positive sign and without normalisation to the modulation sum as an open-loop feedforward, so that $m(\theta)=K_1\sin\theta+K_3\sin^3\theta$. Using the identity $\sin^3\theta=[3\sin\theta-\sin(3\theta)]/4$, the coefficients of $\sin^3\theta$ and of a direct $\sin(3\theta)$ injection are shown to carry different meanings. The Sin\textsuperscript{3}-based feedforward term is added to the modulation reference to reduce the harmonic distortion caused by neutral-switch suppression at low power, and its closed-loop behaviour with current controllers is investigated. A piecewise DCM--CCM average model explains why a sinusoidal duty cycle does not produce a sinusoidal current after suppression. Scanning $0.10\leq k_1,k_3\leq0.50$ in steps of 0.01 with $I_{\mathrm{rms}}>2.2$~A yields 23 coefficient pairs with a third-harmonic ratio below 5\,\%. The best pair, $(k_1,k_3)=(0.36,0.42)$, gives 2.244~A RMS and reduces the analytical third-harmonic ratio from 39.96\,\% for a pure sinusoidal duty cycle to 0.315\,\%, with 4.704\,\% total THD. Normalised to the fundamental, its third-harmonic coefficient is $-0.156$, against $-0.131$ for the calibrated switched-model coefficients. The coefficients depend on the operating point and the model.

\textbf{Closed loop at medium and rated current.} The S4 model uses L1-current feedback with P, PI and PIR controllers in an analogue-equivalent form. Since the S4 model contains no neutral-point balancing loop, a higher DC-bus capacitance is used to limit the midpoint-voltage oscillation. At the rated 25~A, PI and PIR controllers track the reference with $-2.1$\,\% and $-1.0$\,\% amplitude error and 8.5--9.2\,\% THD without saturation, while the P controller leaves $-11.4$\,\%. For a 10~A to 25~A reference step the overshoot relative to the final steady envelope is 0.17\,\% and the settling time 0.43~ms. Neutral-switch suppression is designed for the low-power region.

\textbf{Stability.} The stability analysis shows in three layers that the Sin\textsuperscript{3} term is a known input added from outside the loop and not a controller. Analytically, the characteristic equation $1+L(s)=0$ does not contain the term. In the averaged model and in the linearisation of the switched model, the loop gain $L$, the sensitivity $S$ and the complementary sensitivity $T$ are identical with the term on and off, and in the switched model the differences remain at the level of measurement uncertainty. The effect of the term on the current is set by $S(j\omega)K_{\mathrm{pwm}}G(j\omega)$, and its 150~Hz component is attenuated by $|S|=0.23$. Because the loop gain crosses unity three times (586, 2454 and 3386~Hz), margins are given at every crossing and stability is assessed with the Nyquist criterion and the closed-loop poles. The analogue-equivalent loop at 25~A with PI control has a critical phase margin of 54° at 3.4~kHz, a gain margin of +35~dB at 22~kHz, a delay margin of 44~$\mu$s and a sensitivity peak of 2.4~dB. Its least-damped pole lies at 3.1~kHz with $\zeta=0.18$, and the loop remains stable for 0--5~mH of additional grid inductance.

\textbf{GT-type architecture.} The third contribution re-tasks the T-type switches according to the operating mode, combining the inverter, Vienna PFC and active battery management functions on a common hardware base. The experimental setup, rated 20~kVA with a 20--25~kW design range, is itself built as a GT-type converter: relays and boost diodes in each phase cell select the mode, and multi-secondary battery management transformers serve the charging direction. The low-power experiments were carried out in the inverter (discharge) mode, and the charging chain (PFC, DC bus, transformer primary, secondary and buck/battery) is shown in simulation. Operating modes are defined as safe off, pre-charge, PFC/CC--CV charging, discharge to the grid and relay mode transition. A safe mode transition requires stopped PWM, near-zero phase current, a balanced DC bus and relay position feedback at the same time. By combining the inverter, PFC and battery management functions on shared semiconductors, the architecture can reduce the component count compared with separate power stages.

\textbf{Conclusions.} At low power, the neutral-branch behaviour of the T-type converter offers an efficiency gain at the cost of higher harmonic distortion. The Sin\textsuperscript{3} feedforward reduces this distortion without changing the stability of the current loop. The same T-type semiconductor base can also be organised into a reconfigurable storage architecture. The experiments were carried out under safe laboratory conditions, with the grid side supplied through a variac at about 120~V. The medium- and rated-current results were obtained with the closed-loop model built on the calibrated S3 model.

\textbf{Keywords:} T-type inverter, energy storage system, low-power operation, neutral-switch suppression, feedforward compensation, third harmonic, stability analysis, reconfigurable architecture.
\end{summary}
